\section{MiniCPM：小模型高性能与 muP 稳定 scaling}

为了观察 scaling law 如何真正指导模型训练，本部分先用 MiniCPM 作为完整案例。它把 muP 初始化、learning-rate transfer、optimal batch、WSD schedule 与 Chinchilla-style model/data tradeoff 放进同一条实验链，优点是每个假设都有小规模 sweep 和最终模型结果对应。


\singleslide{slides-images/slide-006.jpg}{MiniCPM：用密集 scaling experiments 与 muP 构造高性能小模型。}{6}
它不是当年最大模型，却公开了 LR、batch、WSD 与 data/model trade-off，适合作为方法案例。
\begin{knowledgebox}{讲义补充：源 Slide 6 为什么选 MiniCPM}
MiniCPM 不是因为最大或最 SOTA 而重要，而是因为它公开展示了小模型如何通过 careful scaling、muP、batch/LR/token-to-size ratio tuning 得到强性能。它给课程提供了“scaling recipes 如何写进真实模型训练”的案例。
\end{knowledgebox}


\singleslide{slides-images/slide-007.jpg}{MiniCPM 的性能定位：1--2.5B 模型可匹配部分 7B，说明 recipe 会移动参数效率曲线。}{7}
Slide 7 的结论不是“小模型永远更好”，而是参数量之外的数据、schedule 与 scaling calibration 会显著改变同规模表现。
\begin{knowledgebox}{讲义补充：源 Slide 7 不能只看模型大小}
这页展示小模型通过训练策略和数据/超参优化达到强性能。它提醒我们：参数量不是性能唯一解释变量。若 scaling recipe 更好、数据更有效、训练 token 更多，小模型可能超过粗糙训练的大模型。
\end{knowledgebox}

\begin{figure}[H]
\centering
\includegraphics[width=0.86\textwidth]{slide-008.jpg}
\caption{Slide 8：MiniCPM 技术 1，使用 muP 稳定 scaling，并给出 scale\_emb、scale\_depth、init\_std、lr 等值。}
\figsource{CS336 Spring 2026 Lecture 11 官方幻灯片。}
\end{figure}

\begin{importantbox}{读图：Slide 8 中 muP 的作用}
muP 的目标是让不同宽度模型的训练动力学更可比，使小模型上找到的 learning rate 和初始化规则更稳定迁移到大模型。这里的具体数值不是通用处方，而是 MiniCPM architecture 和训练 recipe 下的标定结果。
\end{importantbox}


\singleslide{slides-images/slide-009.jpg}{MiniCPM strategy：muP 初始化、固定 aspect ratio，并拟合 batch/LR/token-to-size 后扩大模型。}{9}
最大 sweep 模型与最终模型仍相差约 5 倍，因此成功依赖外推；muP 的作用是减少 width 改变时重新调参。
\begin{knowledgebox}{讲义补充：源 Slide 9 的 strategy 为什么省事}
固定 aspect ratio 意味着不在每个规模都重新搜索 depth/width/head 等比例；muP 则希望 LR 和初始化跨宽度稳定。这样 scaling sweep 的变量更少，外推更容易。但图中也提醒：最大小模型和最终模型之间仍有约 5 倍 gap，外推并非无风险。
\end{knowledgebox}

\subsection{MiniCPM 的 LR、batch 和 WSD}

前面的 muP 只承诺宽度变化时训练动力学更稳定，本节检查这是否真的让最优 learning rate 更可迁移，并进一步拟合 batch 随目标 loss 或数据量的变化。WSD 把 warmup、stable 与 decay 分开，使一次长训练可以在多个中间点比较 model/data 配比，降低从头训练大量候选的成本。

\begin{figure}[H]
\centering
\includegraphics[width=0.86\textwidth]{slide-010.jpg}
\caption{Slide 10：Optimal LR：根据 muP，最优 learning rate 应大致稳定；MiniCPM 检验这一点。}
\figsource{CS336 Spring 2026 Lecture 11 官方幻灯片。}
\end{figure}

\begin{knowledgebox}{读图：Slide 10 应看“最优点是否横向稳定”}
若不同模型规模的 loss-vs-LR 曲线最小点接近同一个 LR，说明 muP 达到了 scale-invariant tuning 的目标。若最优点随规模漂移，说明小模型 LR 外推到大模型仍不可靠。
\end{knowledgebox}

\begin{figure}[H]
\centering
\includegraphics[width=0.86\textwidth]{slide-011.jpg}
\caption{Slide 11：Optimal batch：三个模型大小在数据量、batch 和 loss 上形成曲面，红线标出近似最小 loss。}
\figsource{CS336 Spring 2026 Lecture 11 官方幻灯片。}
\end{figure}

\begin{knowledgebox}{读图：Slide 11 的三维 batch 曲面}
横轴是 batch，纵轴或多条曲线代表数据/训练进度，颜色表示 loss。垂直列是一条固定 batch 的训练曲线。红线追踪不同数据量下的最优 batch。读这页时要看最优 batch 如何随 loss 下降而增大，而不是只找某个全局最小点。
\end{knowledgebox}

\begin{figure}[H]
\centering
\includegraphics[width=0.86\textwidth]{slide-012.jpg}
\caption{Slide 12：Optimal batch size：沿 Kaplan 2020 分析，最优 batch 随最终 loss 降低按多项式增大。}
\figsource{CS336 Spring 2026 Lecture 11 官方幻灯片。}
\end{figure}

\begin{importantbox}{读图：Slide 12 与 critical batch 的关系}
最终 loss 越低，训练越接近后期，critical batch 通常越大。图中的多项式趋势说明 batch size 也可被 scaling law 拟合。它不是说 batch 越大越好，而是说目标 loss 变化时，最优 batch 也应随之 schedule。
\end{importantbox}


\singleslide{slides-images/slide-013.jpg}{剩余问题：模型大小与数据量 trade-off；完整 Chinchilla 网格需要从头训练，成本可从 $n$ 增至 $n^2$。}{13}
早停曲线不能直接代表不同 token budgets 的最终 loss；WSD 尝试复用训练轨迹以减少独立 runs。
\begin{warningbox}{讲义补充：源 Slide 13 的成本爆炸}
若每个模型大小都要从头训练多个数据量点，网格成本会快速变成 \(n_{\text{models}}\times n_{\text{data}}\)。这正是 scaling in practice 的核心痛点：拟合 scaling law 本身也要花 compute，必须寻找复用曲线或降低 sweep 成本的方法。
\end{warningbox}

\begin{figure}[H]
\centering
\includegraphics[width=0.86\textwidth]{slide-014.jpg}
\caption{Slide 14：MiniCPM 的部分解法是 WSD learning rate：warmup、stable、decay 三阶段。}
\figsource{CS336 Spring 2026 Lecture 11 官方幻灯片。}
\end{figure}

\begin{knowledgebox}{术语消化：什么是 WSD learning rate}
WSD 是 warmup-stable-decay。先 warmup 到学习率平台，在 stable phase 长时间训练；需要 Chinchilla-style 不同 token budget 时，可以从 stable 末尾 restart decay，而不是每个数据量都从头训练。它用训练曲线复用降低 scaling sweep 成本。
\end{knowledgebox}

\begin{figure}[H]
\centering
\includegraphics[width=0.86\textwidth]{slide-015.jpg}
\caption{Slide 15：WSD 在 MiniCPM 中表现良好，stable 阶段慢，decay 阶段 loss 快速下降，decay 约 10\%。}
\figsource{CS336 Spring 2026 Lecture 11 官方幻灯片。}
\end{figure}

\begin{knowledgebox}{读图：Slide 15 的曲线说明什么}
WSD 的好处不是让 stable 阶段最快，而是让同一条训练轨迹在多个候选 token budget 上可复用。Decayed tail 给出类似“如果在这里结束训练”的 loss 点，从而支持 lower-envelope 或 joint fit。
\end{knowledgebox}

\subsection{MiniCPM 的 Chinchilla 分析}


\singleslide{slides-images/slide-016.jpg}{MiniCPM 的 Chinchilla-type analysis：借助 WSD 使用 lower envelope 与 joint fit。}{16}
两种方法分别从观测最优点和全网格函数出发；一致性比单个拟合常数更重要。
\begin{knowledgebox}{讲义补充：源 Slide 16 的方法选择}
Method 1 是 lower envelope：看所有训练曲线在不同 compute 下的最小点。Method 3 是 joint fit：拟合参数量和数据量到 loss 的曲面。MiniCPM 同时使用两者，是为了在降低成本的同时增强外推可信度。
\end{knowledgebox}


\singleslide{slides-images/slide-017.jpg}{Chinchilla Method 1：不同模型 training curves 的 lower envelope 与数据边际收益。}{17}
颜色代表模型规模；envelope 是否直线决定 power-law 外推是否可信，图中轻微曲率应保留为不确定性。
\begin{knowledgebox}{讲义补充：源 Slide 17 的颜色和趋势}
不同颜色代表不同模型规模。若 lower envelope 呈直线或平滑曲线，就可以外推 compute-optimal 关系。图中趋势清晰但不完美，说明公开 scaling recipe 在真实数据中会有噪声，不能只看拟合线的漂亮程度。
\end{knowledgebox}

\begin{figure}[H]
\centering
\includegraphics[width=0.86\textwidth]{slide-018.jpg}
\caption{Slide 18：Chinchilla method 3，MiniCPM 的主要 scaling approach 是 joint fit，并得到很高 data-model ratios。}
\figsource{CS336 Spring 2026 Lecture 11 官方幻灯片。}
\end{figure}

\begin{importantbox}{读图：Slide 18 的高 data-model ratio 意味着什么}
如果 joint fit 推荐很高 token-to-parameter ratio，训练策略会偏向“小一些模型 + 多训练数据”。这与 Chinchilla 以来的过训趋势一致，也解释为什么小模型通过更多 token 和精细 tuning 能表现很强。
\end{importantbox}

